\documentclass[12pt, a4paper]{article} % Classe: article, report o book [8]

% --- Lingua e Codifica (Supporto Italiano) ---
\usepackage[utf8]{inputenc}  % Codifica input
\usepackage[T1]{fontenc}      % Codifica font
\usepackage[italian]{babel}   % Lingua italiana [2]

% --- Impaginazione e Margini ---
\usepackage[margin=2.5cm]{geometry} % Imposta margini (2.5cm su tutti i lati)

% --- Pacchetti Scientifici/Matematici ---
\usepackage{amsmath, amssymb, amsfonts, amsthm} % Matematica avanzata [2]
\usepackage{graphicx} % Per inserire immagini
\usepackage{hyperref} % Per collegamenti ipertestuali
\usepackage{booktabs} % Per tabelle professionali

% --- Formattazione Testo e Altro ---
\usepackage[document]{ragged2e}
\usepackage{xcolor}    % Per usare i colori 


\title{Isacco Barbieri}
\date{27-01-2026}


\begin{document}

\maketitle

\section{Ruolo}
Isacco Barbieri è un Agente esterno con esperienza di oltre 10 anni di
collaborazione con l'azienda e si occupa di :
\begin{itemize}
      \item Creazione delle offerte Access per conto di clienti della sua zona di
            competenza.
      \item Gestione delle relazioni con i clienti.
      \item Supporto post-vendita e assistenza ai clienti.
\end{itemize}

\section{Situazione attuale}
\subsection{Attuale uso di Access}
Isacco utilizza Access per:
\begin{itemize}
      \item Creare le offerte per i clienti.
      \item Creazione delle offerte di assistenza e ampliamenti.
      \item Aggiunta clienti e gestione anagrafiche.
      \item Trasfomazione delle offerte in ordini se accettate.
\end{itemize}

\subsection{Attuale vista dati in Access}
Isacco all'interno dell'attuale programma ha la possibilità di vedere:
\begin{itemize}
      \item Le offerte create da lui.
      \item Le richieste di assistenza (RA) create da lui.
      \item L'anagrafica clienti completa.
      \item L'elenco articoli completo.
      \item Le offerte e RA degli altri agenti non sono visibili.
\end{itemize}

\subsection{Gestione utenti e visibilità}
Per quanto riguarda la gestione degli utenti e la visibilità dei dati, Isacco
ha evidenziato i seguenti punti chiave:
\begin{itemize}
      \item \textbf{Segregazione dei dati:} È fondamentale mantenere la distinzione tra offerte e RA (Richieste assistenza/ricambi).
            La distinzione tra i due tipi aiuta a evitare confusione e garantisce che le operazioni siano gestite correttamente.
      \item \textbf{Privacy per autore:} Un agente esterno non deve vedere i progetti o le offerte degli altri agenti per evitare conflitti di interesse (Policy aziendale).
      \item \textbf{Anagrafica clienti condivisa:} A differenza delle offerte, l'elenco clienti deve essere visibile a tutti,
            può capitare di fare offerte per conto di altri commerciali.
\end{itemize}

\section{Requisiti emersi}
\subsection{Accessibilità e modalità offline}
È stato sollevato il problema legato alla connessione. Attualmente i commerciali esterni si collegano ad un computer interno tramite VPN per utilizzare Access.
Tuttavia, in caso di perdita di connessione (ad esempio in mobilità), la connessione cade e serve ripristinarla . Il lavoro non viene perduto, ma l'operatività viene interrotta.
è stato richiesto che :
\begin{itemize}
      \item Il nuovo sistema dovrebbe permettere di lavorare anche in assenza di
            connessione Internet per una vendita o un'assistenza urgente.
      \item Una volta ristabilita la connessione, il sistema dovrebbe sincronizzare
            automaticamente i dati senza richiedere interventi manuali.
\end{itemize}

\subsection{Gestione articoli e codifica}
Per quanto riguarda la parte di creazione offerte è emerso che:
\begin{itemize}
      \item \textbf{Codici generici:} Necessità di voci generiche per categorie (es. "Valvola generica", "Vibratore generico") per gestire articoli su misura
            senza dover creare migliaia di codici specifici che "sporcano" il database. Questi articoli devono essere liberamente prezzati.
      \item \textbf{Articoli "Ex" (Obsoleti):} Le voci obsolete non devono essere eliminate, in quanto potrebbero servire per la manutenzione di vecchi impianti o richieste
            specifiche dei clienti. Devono essere però facilmente individuate per non incorre in errori.
      \item \textbf{Clienti:} La gestione della creazione dei cliente deve essere ponderata per evitare duplicati o errori di inserimento.
\end{itemize}

\subsection{Workflow operativo: offerte e ordini}
Sul flusso di lavoro operativo, Isacco ha evidenziato:
\begin{itemize}
      \item \textbf{Copia e template:}
            \begin{itemize}
                  \item \textbf{Funzionalità essenziale:} Possibilità di copiare interamente vecchie offerte o
                        RA per risparmiare tempo su progetti simili.
                  \item \textbf{Ipotesi template:} L'uso di template per la creazione o suggerimento componenti è benvenuto, ma la copia dello storico rimane prioritaria.
            \end{itemize}
      \item \textbf{Revisioni:} Necessità di gestire versioni multiple della stessa offerta (Revisione 1, 2, ecc.) mantenendo lo storico delle modifiche senza sovrascrivere l'originale.
            In fase contrattuale capita frequentemnete di dover fare più offerte con lievi modifiche. Le varie versioni/revisioni devono matenere sempre lo stesso numero di offerta.
      \item \textbf{Conferma commesse:} Gli esterni possono confermare le RA, ma per le commesse complesse rimane il vincolo della gestione fisica delle cartelle interne \textit{(Problema interno organizzativo)} per cui non possono direttamnete confermarle,
            ma solo attraverso un commerciale interno. Questa procedure fa tempo e si accorcerebbero i tempi con un sistema diverso.
      \item \textbf{Creazione Offerta per terzi:} Necessità di creare offerte per conto di altri commerciali in caso di assenza o supporto, mantenendo la tracciabilità dell'autore originale.
            Questa possibilità deve esserci per tutti i tipi di commerciali sia interni che esterni.
      \item \textbf{Automazione conferma RA:}
            Attualmente il processo è manuale e frammentato \textit{(stampa PDF,
                  salvataggio in cartelle, invio email manuale).}:
            \begin{itemize}
                  \item \textbf{Richiesta:} Un sistema automatizzato che, alla conferma dell'ordine, invii notifiche e documenti ai reparti competenti
                        \textit{(Amministrazione, Ordini/Benazzi)} con la possibilità di allegare note e file.
            \end{itemize}
      \item \textbf{Infomazioni tecniche:} I clienti in fase di offerta richiedono a volte le potenze e il consumo di aria ,
            serve manternere aggiurnate le infomazioni tecniche per dare la possibilità al commerciale di quantificarle  preliminarmente in questa fase.
\end{itemize}

\subsection{Gestione tecnica post-vendita}
Alla proposta di una gestione dei contratti revisionati dai tecnici, isacco ha
accolto positivamente l'idea:
\paragraph{Contesto Attuale:}
Attualmente la rilevazione delle modifiche tecniche non vengono fatte dai
commerciali interni , ma quelli esterni richiedono espressamnete che venga
valutata dal tecnico che invia tramite e-mail il materiale. Dopo la valutazione
tecnica si decide se emettere una nuova offerta o preventivo di modifica.
\paragraph{Richiesta:}
\begin{itemize}
      \item \textbf{Integrazione modifiche tecniche:} L'ufficio tecnico (PE) deve poter integrare le modifiche tecniche direttamente nel sistema gestionale,
            creando una "copia tecnica" parallela al contratto di vendita.
      \item \textbf{Aggiornamento offerta/contratto:} Il sistema deve pertmettere di confrontare facilmente la versione originale con quella modificata,
            evidenziando le differenze e le aggiunte. Dato il confronto deve dare la possibilità al commerciale di generare una nuova offerta basata sulle
            modifiche tecniche apportate come revisione dell'originale o come richeista di assistenza.
      \item \textbf{Alert commerciale:} Al completamnento del contratto tecnico il commerciale di riferimento deve essere notificato per valutare le modifiche.
\end{itemize}
Questo sistema aiuterebbe a ridurre errori e discrepanze tra la vendita e la realizzazione portando vantaggi come:
\begin{itemize}
      \item Maggiore accuratezza nelle offerte e nei contratti successivi avendo degli
            esempi corretti a livello tecnico.
      \item Riduzione delle controversie con i clienti dovute a incomprensioni tecniche.
      \item Migliore tracciabilità delle modifiche e delle decisioni prese durante il ciclo
            di vita del contratto.
      \item Margini di profitto più chiari e definiti grazie a una migliore gestione delle
            variazioni tecniche.
\end{itemize}

\subsection{Integrazione strumenti esterni (Excel)}
Ad oggi per il calcolo dei costi di determinati componenti (es. Silos, Scale,
Quadri Elettrici) si utilizzano complessi fogli di calcolo Excel sviluppati
dagli specialisti di prodotto. Isacco suggerisce che:
\begin{itemize}
      \item \textbf{Mantenimento File Excel:} I fogli di calcolo attuali (per Silos, Scale, Quadri Elettrici) sono molto complessi,
            contengono logiche di verifica e vengono aggiornati frequentemente dagli specialisti.
      \item \textbf{Decisione Progettuale:} Non conviene tentare di replicare questa logica complessa nel nuovo gestionale (costoso e rischioso).
            È preferibile continuare a usare gli Excel e inserire nel gestionale solo il risultato finale (prezzo/descrizione).
      \item \textbf{Vantaggi dei fogli attuali:}
            I fogli di calcolo attuali hanno il vantaggio di:
            \begin{itemize}
                  \item \textbf{Flessibili:} Possono essere modificati rapidamente dagli specialisti senza dover coinvolgere il reparto IT.
                  \item \textbf{Specializzati:} Contengono conoscenze specifiche di prodotto che sarebbero difficili da trasferire in un sistema gestionale generico.
                  \item \textbf{Blocco Errori:} Hanno logiche di controllo e verifica che impediscono di vendere configurazioni non valide.
            \end{itemize}
\end{itemize}
L'internalizzazione di questi calcolo all'interno del configuratore è stata suggerita di intraprenderla in una fase successiva, una volta che il sistema base sia
stabile e funzionante.

\subsection{Gestione sconti e ricarichi}
Dare la possibilità di gestire sconti e ricarichi in modo flessibile:
\begin{itemize}
      \item Necessità di gestire sconti visibili al cliente e ricarichi "nascosti" (per
            coprire costi extra).
      \item Possibiltà di fare sconti solo su materiale e sconti comrpensivi di manodopera
            di montaggio
      \item Possibilità di imputare sconti al singolo articolo o all'intera offerta.
      \item Attualmente è possibile imputare un prezzo concordato all'intera offerta e
            questu'ultima si calcola gli sconti in modo automatico. Funzione da mantenere.
      \item Possibilità di scontare anche i singoli capitoli dell'offerta.
\end{itemize}

\subsection{Condizioni di vendita predefinite}
Alcuni clienti hanno delle condizioni di vendita fisse (es. sconto del 10\% su
tutto il materiale) che possono essere suggerite automaticamente in fase di
creazione offerta.
\begin{itemize}
      \item Possibilità di associare condizioni di vendita predefinite a clienti specifici
            (es. Sconto del 10\% su tutto il materiale per il cliente X).
      \item Durante la creazione dell'offerta, il sistema dovrebbe suggerire
            automaticamente queste condizioni in base al cliente selezionato.
      \item L'agente deve avere la possibilità di modificare o rimuovere queste condizioni
            se necessario, ma con un avviso che indica che si sta discostando dalle
            condizioni standard del cliente.
\end{itemize}

\subsection{Stampa offerte e layout personalizzati}
\begin{itemize}
      \item \textbf{Capitoli:} I clienti richiedono offerte con prezzi dettagliati per "Capitolo" \textit{(es. Totale Linea Silos)}, nascondendo i prezzi dei singoli componenti.
            La scelta dei capitoli deve essere flessibile e personalizzabile.
      \item \textbf{Layout di Stampa:} Flessibilità nel mostrare/nascondere descrizioni tecniche interne, prezzi unitari o totali a seconda della tipologia di documento (Offerta vs Lista Ricambi).
\end{itemize}

\end{document}